Existing service invocation, publication, and access-control mechanisms provide the basis for NDNSF's candidate discovery and permission management.

RICE~\cite{krol2018rice} studies remote method invocation in ICN, including an invocation handshake and parameter/result handling. NDNSF builds on these established ideas for invocation and protected parameter exchange.

DNMP~\cite{nichols2019dnmp} publishes measurement commands through Sync and collects replies. VerSec relates allowed command components to role signing keys. The schema determines which roles may issue commands with particular targets and parameters.

NDNSF instead uses service-attribute encryption to restrict descriptor readers and aggregates permissions in DKEYs. Adding encryption to a DNMP-style design would also protect command content, with additional key-distribution costs.

NSC~\cite{meirovitch2021nsc} requires a signed initial invocation Interest, uses application-defined trust schemas, and describes encrypted input/result retrieval with NAC support. Its report also discusses retry behavior for long-running computation.

NDNSF explicitly collects candidate feedback before selection. This differs from the basic NSC invocation sequence, but applications can extend an invocation system with discovery or negotiation; their complete costs must be compared.

NDNSD~\cite{dulal2022ndnsd} supports service advertisement discovery and optional NAC-ABE protection of service information. NDNSF can use these advertisements for provider preselection, while its Request--ACK exchange obtains participation feedback for a particular task.

NAC~\cite{zhang2018nac} automates named-data access-key distribution and analyzes ABE-based aggregation. Enhanced NAC-ABE supplies KP-ABE, policy-bearing DKEYs, validation support, and key segmentation~\cite{dulal2023nac-abe}.

NDNSF applies these capabilities to service descriptors and invocation permissions. The reduction in credential objects follows from ABE permission aggregation. Section~\ref{sec:credential-cost} examines this benefit and its key-distribution and update costs.

Trust schemas can encode general rules rather than one rule per entity or service~\cite{yu2015trust}. Other designs can also reuse identity credentials and general trust rules. NDNSF keeps service-permission lists in ABE policies rather than baseline trust rules; this does not establish lower rule-management or validation costs than every alternative.
